\documentclass[11pt]{article}
\usepackage[margin=1in]{geometry}
\usepackage{amsmath,amssymb,amsthm,hyperref}
\hypersetup{colorlinks=true,urlcolor=blue}
\newtheorem{theorem}{Theorem}
\newtheorem{corollary}[theorem]{Corollary}
\newcommand{\dd}{\,dt}
\title{A gamma quotient with strict logarithmic complete monotonicity}
\author{Research record: MathDB 336698}
\date{12 September 2026}
\begin{document}
\maketitle

\section*{Task 1. Resolve the gamma-quotient conjecture}
Arora and Sahoo, \emph{Interpolation on Gauss hypergeometric functions
with an application}, Involve \textbf{11}(4) (2018), 625--641,
\href{https://doi.org/10.2140/involve.2018.11.625}{doi:10.2140/involve.2018.11.625},
pose Conjecture 3.14 on page 635: for every $0\le a\le1$, prove that
\[
 Q_a(c)=\frac{\Gamma(c)\Gamma(c-4)}
 {\Gamma(c-a)\Gamma(c+a-1)}
\]
decreases on $c>4$. The answer is affirmative, and strict.

\begin{theorem}\label{thm:main}
For every $-3\le a\le4$, $c>4$, and integer $m\ge1$,
\[
 (-1)^m\frac{d^m}{dc^m}\log Q_a(c)>0.
\]
In particular, $Q_a'(c)<0$ throughout the conjectured domain, including
the parameter endpoints $a=0,1$.
\end{theorem}

\begin{proof}
Write $\psi=\Gamma'/\Gamma$. The classical digamma difference formula is
\[
 \psi(u)-\psi(v)
 =\int_0^\infty
 \frac{e^{-vt}-e^{-ut}}{1-e^{-t}}\dd,\qquad u,v>0.
\]
For example, it follows by subtraction from DLMF
\href{https://dlmf.nist.gov/5.9.E16}{5.9.16}; equivalently,
expand $(1-e^{-t})^{-1}$ and integrate the usual digamma series.
All four gamma arguments here are positive. Consequently
\begin{equation}\label{eq:kernel}
 -\frac{d}{dc}\log Q_a(c)
 =\int_0^\infty e^{-ct}
 \frac{1+e^{4t}-e^{at}-e^{(1-a)t}}{1-e^{-t}}\dd.
\end{equation}
Fix $t>0$. The function
$a\mapsto e^{at}+e^{(1-a)t}$ is convex and takes the same value
$e^{-3t}+e^{4t}$ at the endpoints of $[-3,4]$. It is therefore at most
that value throughout the interval. Thus
\[
 1+e^{4t}-e^{at}-e^{(1-a)t}
 \ge 1-e^{-3t}>0.
\]
At $t=0$, the numerator in \eqref{eq:kernel} is $3t+O(t^2)$,
so the quotient is bounded. At infinity it is $O(e^{4t})$.
For $c>4$, the integral therefore converges absolutely, and it may be
differentiated any number of times, locally uniformly in $c$, by
exponential domination. For every $m\ge1$ this gives
\[
 (-1)^m(\log Q_a)^{(m)}(c)
 =\int_0^\infty t^{m-1}e^{-ct}
 \frac{1+e^{4t}-e^{at}-e^{(1-a)t}}{1-e^{-t}}\dd>0.
\]
Finally $Q_a>0$, so $Q_a'=Q_a(\log Q_a)'<0$.
\end{proof}

\paragraph{A short independent proof on the requested range.}
The classical series
$\psi'(x)=\sum_{k=0}^\infty(x+k)^{-2}$ shows that $\psi$ is increasing
and strictly concave on $(0,\infty)$
(DLMF \href{https://dlmf.nist.gov/5.15.E1}{5.15.1}).
For $0\le a\le1$, the concave function
$a\mapsto\psi(c-a)+\psi(c+a-1)$ has equal endpoint values.
Hence
\[
 \psi(c-a)+\psi(c+a-1)\ge\psi(c)+\psi(c-1),
\]
and therefore
\[
 (\log Q_a)'(c)
 \le\psi(c-4)-\psi(c-1)<0.
\]
This also proves the entire original conjecture directly.

\begin{corollary}[Sharp parameter interval]
On the domain where all four gamma arguments are positive,
\[
 c>b_a:=\max\{4,a,1-a\},
\]
$Q_a$ is decreasing everywhere if and only if $-3\le a\le4$.
For precisely these parameters it is strictly logarithmically
completely monotone.
\end{corollary}
\begin{proof}
Sufficiency is Theorem \ref{thm:main}. If $a>4$, then $b_a=a$ and
$Q_a(c)\to0$ as $c\downarrow a$, because $\Gamma(c-a)\to\infty$
while the other three factors have finite positive limits.
A positive decreasing function cannot have left-endpoint limit zero.
If $a<-3$, the same argument at $c=1-a$ uses
$\Gamma(c+a-1)\to\infty$. Thus global decrease fails outside the
stated interval, and so does logarithmic complete monotonicity.
\end{proof}

\paragraph{Status and scope.}
The published source retains the conjecture, and the live MathDB
record was open with no solution or discussion entries at the dated
checks recorded in the accompanying literature audit. No matching
full public claim was located in that bounded search.
The proof uses classical digamma formulas and elementary convexity;
it does not claim those general tools as new.

\section*{Task 2. Find a positive inverse Laplace representation}
For each $a\in[-3,4]$, derive an explicit positive measure representing
$Q_a(4+x)$ as a Laplace transform. Determine how the measure's density
and endpoint behavior depend on $a$.

\section*{Task 3. Determine generalized Stieltjes orders}
For each $a\in[-3,4]$, determine the least $\alpha>0$, if one exists,
for which $Q_a(4+x)$ has a representation
$\int_0^\infty(x+t)^{-\alpha}\,d\mu_a(t)$ with $\mu_a\ge0$.
Include the rational endpoint cases as exact benchmarks.

\section*{Task 4. Sharpen the interpolation error bounds}
Use the positive kernel to obtain optimal uniform bounds for
$-(\log Q_a)'(c)$ on $0\le a\le1$, and propagate these bounds into
the quadratic hypergeometric interpolation errors in the source.

\section*{Task 5. Compare decrease with logarithmic complete monotonicity}
For quotients of three gamma factors by three gamma factors with fixed
total shift difference, find explicit parameter families that are
decreasing on their positive-argument domain but fail logarithmic
complete monotonicity, or prove their equivalence on a specified family.
\end{document}
